\documentclass[11pt,a4paper,oneside,notitlepage]{report}
% (преамбула винесена з main.tex; спільна для main і драйверів розділів)

% ---------- мова і шрифти ----------
\usepackage{fontspec}
\usepackage{polyglossia}
\setdefaultlanguage{ukrainian}
\setotherlanguage{english}
\gappto\captionsukrainian{\renewcommand{\bibname}{Бібліографія}\renewcommand{\refname}{Бібліографія}}
\setmainfont{cmunrm}[Extension=.otf,BoldFont=cmunbx,ItalicFont=cmunti,BoldItalicFont=cmunbi]
\setsansfont{cmunss}[Extension=.otf,BoldFont=cmunsx,ItalicFont=cmunsi,BoldItalicFont=cmunso]
\setmonofont{cmuntt}[Extension=.otf,BoldFont=cmuntb,ItalicFont=cmunit,BoldItalicFont=cmuntx]

% ---------- верстка ----------
\usepackage[a4paper,top=2cm,bottom=2.2cm,left=2.2cm,right=2cm]{geometry}
\usepackage{titlesec}
\titleformat{\chapter}[hang]{\Large\bfseries\sffamily}{\thechapter.}{0.6em}{}
\titlespacing*{\chapter}{0pt}{-10pt}{12pt}
\titleformat{\section}{\large\bfseries\sffamily}{\thesection}{0.6em}{}
\titleformat{\subsection}{\normalsize\bfseries\sffamily}{\thesubsection}{0.6em}{}
\usepackage{enumitem}
\usepackage{tocloft}
\usepackage{multicol}
\setlength{\cftbeforechapskip}{2pt}
\renewcommand{\cftsecfont}{\small}\renewcommand{\cftsecpagefont}{\small}
\setlength{\cftbeforetoctitleskip}{0pt}
\renewcommand{\cfttoctitlefont}{\large\bfseries\sffamily}
\setlength{\cftaftertoctitleskip}{4pt}
\setlength{\cftbeforesecskip}{0pt}

\setlist{itemsep=1pt,topsep=3pt,parsep=0pt}
\setlength{\emergencystretch}{3em}% довгі шляхи файлів у \texttt
\usepackage{booktabs,array,tabularx,longtable}
\renewcommand{\arraystretch}{1.12}
\usepackage{graphicx}
\graphicspath{{figures/}}
\usepackage[font=small,labelfont=bf,labelsep=period]{caption}
\usepackage{subcaption}
\usepackage[table]{xcolor}
\usepackage{tikz}
\usetikzlibrary{arrows.meta,calc,decorations.markings,positioning,shapes.geometric}

% ---------- математика і позначення ----------
\usepackage{notation}
\numberwithin{equation}{chapter}

% ---------- рамки ----------
\usepackage[most]{tcolorbox}
\tcbset{enhanced,breakable,boxrule=0.5pt,arc=1.5pt,left=5pt,right=5pt,top=3pt,bottom=3pt,
        fonttitle=\sffamily\bfseries\small\hypersetup{citecolor=white,linkcolor=white},before skip=6pt,after skip=6pt}
% стандартна теорія поза корпусом (рішення Q1)
\newtcolorbox{outofcorpus}[1][]{colback=gray!5,colframe=gray!60!black,
  title={Поза корпусом: стандартна теорія\ifx&#1&\else{} --- #1\fi}}
% приклад з корпусу: реальні установки й числа
\newtcolorbox{corpusexample}[1][]{colback=blue!3,colframe=blue!50!black,
  title={Приклад з корпусу\ifx&#1&\else{} --- #1\fi}}
% виведення з пропущеними рутинними кроками
\newtcolorbox{derivation}[1][]{colback=white,colframe=black!40,
  title={Виведення\ifx&#1&\else{} --- #1\fi}}
% практикум на коді проєкту
\newtcolorbox{practicum}[1][]{colback=green!4,colframe=green!40!black,
  title={Практикум\ifx&#1&\else{} --- #1\fi}}
% дослідницький трек: статуси з реєстрів Our try/00-decisions
\newtcolorbox{established}[1]{colback=teal!4,colframe=teal!60!black,
  title={Встановлено в проєкті [#1]}}
\newtcolorbox{conjecture}[1]{colback=orange!6,colframe=orange!70!black,
  title={Гіпотеза [#1] --- не перевірено}}
\newtcolorbox{refuted}[1]{colback=red!4,colframe=red!55!black,
  title={Спростовано [#1] --- повчальний приклад}}


% ---------- код (путівник) ----------
% minted (Pygments через latexminted); збірка: latexmk (див. latexmkrc: -shell-escape, PATH до Pygments).
% listings відмовились: під XeLaTeX він зсуває нумерацію, якщо над фрагментом є не-ASCII символи.
\usepackage{minted}
\usepackage{needspace}
\setminted{fontsize=\scriptsize,linenos,numbersep=6pt,breaklines,breakanywhere,frame=single,
  framesep=2mm,rulecolor=black!25,bgcolor=black!2,style=default}
% \code{шлях}{перший}{останній}{підпис} — фрагмент БЕЗПОСЕРЕДНЬО з файлу проєкту (через src/-посилання)
\newcommand{\code}[4]{\par\needspace{7\baselineskip}\noindent
  {\footnotesize\texttt{\detokenize{#1}}, рядки #2--#3\quad #4}\par\nopagebreak[4]\vspace{-2pt}
  \inputminted[firstline=#2,lastline=#3,firstnumber=#2]{python}{src/#1}}
% ---------- картки пунктів реєстру ----------
\newtcolorbox{itemcard}[2]{colback=black!2,colframe=black!55,title={#1\quad\normalfont\small #2}}

% ---------- посилання ----------
\usepackage[colorlinks=true,linkcolor=blue!50!black,citecolor=green!40!black,
            urlcolor=blue!60!black,bookmarksnumbered]{hyperref}
\usepackage{xurl}
\usepackage[ukrainian,nameinlink]{cleveref}

\title{\sffamily\bfseries Фізика плазми всередині токамака\\[4pt]
       \large Навчальний посібник (методичка)}
\author{На основі корпусу джерел проєкту «Моделювання потоку плазми в контурі токамаку»}
\date{2026}

